\section{Introduction}

\begin{itemize}

  \item Privacy-preserving ledgers still need to support auditing: regulators,
    central banks, and companies all need to recover specific transactions on
    demand.

  \item Existing approaches sit at two extremes: fully transparent ledgers give
    up all privacy, while a long-term view key gives an auditor unconditional,
    permanent visibility and becomes a single point of failure.

  \item Real audits are narrower: they target one user during one period, so
    authorisation to trace should not outlive that period, and should not
    require the user's continued involvement.

  \item We want a middle ground: per-epoch, self-service tracing authority that
    expires automatically.

  \item Each user holds a long-term tracing key, certified once, from which
    they derive a fresh epoch tracing key every epoch.

  \item Transactions are tagged with a PRF over the epoch key; disclosing that
    key lets an auditor trace exactly the transactions from one epoch, and
    nothing else.

  \item Granting tracing authority is a one-time key disclosure to the chosen
    auditor; no further interaction with the registration authority, ledger
    operator, or user is needed, and authorisation lapses automatically when
    the epoch ends.

  \item We call this mechanism tracing rather than collection: the auditor only
    ever learns one epoch's transactions, never a user's full history. This
    matches existing terminology, where tracing means recovering every
    transaction linked to a user, as distinct from the narrower anonymity
    revocation, which recovers only the identity behind a single transaction.

\end{itemize}

\subsection{Related work}\label{sec:txid-related}

Prior work on regulatory oversight in privacy-preserving payment systems
typically supports either anonymity revocation or tracing.  Anonymity
revocation lets an auditor identify the user responsible for a specific
transaction; tracing, by contrast, lets the auditor recover all transactions
associated with a given user. Tracing can be emulated by applying anonymity
revocation to every transaction in turn, but this costs time linear in the
number of transactions and exposes information about non-targeted users.

\paragraph{Anonymity revocation.} Many protocols attach a verifiable encryption
or other hidden user identifier to each transaction, so that an auditor holding
the corresponding secret key can selectively recover the sender's
identity~\cite{SCN:CamHohLys06,CANS:BDET21,INDOCRYPT:CDLP22}.  Tracing a user's
activity this way requires scanning and decrypting every transaction on the
ledger, incurring overhead linear in the ledger size and exposing unrelated
users' activity to the auditor. These protocols enable identity recovery, but
not efficient targeted tracing.

\paragraph{Tracing via sequential tags.} Other protocols derive per-transaction
tags from a user's tracing secret combined with a monotonically increasing
counter, and use zero-knowledge proofs to show the tags are correctly formed
and bound to the transaction's origin~\cite{CCS:KiaKohSar22,CANS:SarKiaKoh24}.
Given the tracing secret of a target user, an auditor identifies the
corresponding transactions by recomputing the tags and matching them against
the ledger, and the counter binding also stops a user later repudiating their
own transactions. Similar tag-based techniques trace asset lineage rather than
user activity~\cite{FC:GarGreMie16}. All of these share the same limitation: a
user's transactions must be generated sequentially, which prevents concurrent
transaction creation and degrades performance in high-throughput settings.

\paragraph{Auditable privacy systems.} Other protocols provide auditability
through interactive protocols or global consistency checks. ZKLedger lets
auditors verify transaction correctness, but requires interaction and incurs
overhead linear in the number of participants~\cite{EPRINT:NarVasVir18}.
Privacy-focused cryptocurrencies such as Monero, Zcash, and Quisquis achieve
strong anonymity guarantees but do not directly support efficient
tracing~\cite{EPRINT:Noether15,SP:BCGGMT14,AC:FMMO19,FC:BAZB20}.  Retrofitting
tracing into such systems typically requires additional interaction, user
cooperation, or linear-time analysis~\cite{ACNS:PBGP25}. Platypus and ZKAT
incorporate regulatory policy enforcement, for instance balance limits, but
neither provides a mechanism for efficient per-user tracing across arbitrary
transactions~\cite{CCS:WKDC22,EPRINT:ACDDET19}.

\paragraph{Our contribution.} We introduce a tracing mechanism that eliminates
the need for encrypted user identifiers and sequential counters, supporting
efficient tracing while letting users generate and submit multiple transactions
in parallel without restriction. The design gives users fine-grained control
through epoch-bound tracing permissions: these let an auditor identify every
transaction linked to a specific user within a predefined time window, and
automatically become invalid once that epoch expires, requiring a fresh
permission to be issued for any subsequent one. We formalise tracing as an
ideal functionality capturing anonymity, non-repudiation, and permission
expiration, give a black-box construction from signatures, commitments,
pseudorandom functions, and zero-knowledge proofs, and instantiate it
concretely in pairing-based groups with a transparent setup, using the BBS+
signature scheme to let the registration authority certify a user's keys
without ever learning their tracing key.

\paragraph{Layout.} Section~\ref{sec:txid-definitions} formalises the problem
as an ideal functionality. Section~\ref{sec:txid-construction} gives a
construction from generic building blocks.
Section~\ref{sec:txid-instantiation} gives a concrete instantiation in
pairing-based groups, proved secure in Section~\ref{sec:txid-security}.
Section~\ref{sec:txid-evaluation} evaluates the concrete instantiation.
